%%%%%%%%%%%%%%%%%%%%%%%%%%%%%%%%%%%%%%%%%%%%%%%%%%%%%%%%%%%%%%
%%%%%%%%%%%%%%%%%% MA-0838 Geometría Diferencial %%%%%%%%%%%%%%%%%%%%%%%%%%%%%%
%%%%%%%%%%%%%%%%%%%%%%%%%%%%%%%%%%%%%%%%%%%%%%%%%%%%%%%%%%%%%%
% Combina la Opción 2 (Chaves) con el programa MA-0870 (carta al
% estudiantado, II ciclo 2026): variedades, formas, integración y
% geometría riemanniana.
\newpage
\subsection*{MA-0838 Geometría Diferencial}
\addcontentsline{toc}{subsection}{MA-0838 Geometría Diferencial}

\hypertarget{MA0838}{}
\begin{refsection}
\begin{table}[H]
\centering{
\begin{tabular}{|ll|}
\hline
\textbf{Año:} IV  & \textbf{Requisitos:} MA-0471 Introducción a la Geometría Diferencial, \\
 & MA-0556 Ecuaciones Diferenciales y MA-0635 Topología General\\
\textbf{Ciclo:} Optativa  & \textbf{Correquisitos:} Ninguno \\
\textbf{Tipo de curso:} Teórico & \textbf{Horas presenciales por semana:} 4 \\
\textbf{Créditos:} 4 & \textbf{Horas de trabajo independiente por semana:} 8 \\
\textbf{Virtualidad:} P  & \textbf{Clasificación:} Curso propio \\ \hline
\end{tabular}
}
\end{table}

\subsubsection*{Descripción del curso}

La geometría diferencial estudia espacios geométricos localmente euclidianos dotados de una estructura suave, lo cual permite trasladar a ellos las herramientas del cálculo en varias variables. Su objeto central es la \emph{variedad diferencial}, que proporciona el marco natural para generalizar las curvas y superficies a espacios de dimensión arbitraria y para estudiar sus propiedades de manera intrínseca, independientemente de un espacio ambiente.

Este curso optativo se ubica en el cuarto año, una vez que la persona estudiante dispone de la geometría local de curvas y superficies (MA-0471) y de las herramientas de ecuaciones diferenciales ordinarias (MA-0556). Introduce herramientas fundamentales de la disciplina: espacios tangentes, campos vectoriales, fibrados vectoriales, formas diferenciales, métricas riemannianas, conexiones, curvatura e integración. Se desarrollan resultados representativos como el teorema del valor regular, el lema de Poincaré y el teorema de Stokes, y se retoma el teorema egregio en relación con la curvatura riemanniana.

La geometría diferencial desempeña un papel central en áreas como la geometría riemanniana y simpléctica, la topología diferencial, los grupos de Lie y los sistemas dinámicos. En la física, proporciona el lenguaje natural para formulaciones geométricas de la mecánica clásica, la relatividad general y teorías gauge; sus métodos también encuentran aplicaciones en robótica y teoría de control.

El énfasis es teórico: definiciones precisas, enunciados completos y demostraciones de los resultados fundamentales indicados en la metodología. Se espera que la persona estudiante formule argumentos rigurosos, conecte la geometría intrínseca con el cálculo diferencial en variedades y reconozca el alcance de cada hipótesis.

\subsubsection*{Objetivos}

\textbf{General:} Comprender y aplicar las herramientas fundamentales del cálculo diferencial e integral y de la geometría riemanniana en variedades, demostrando resultados centrales.

\textbf{Específicos:} Al finalizar el curso se espera que la persona estudiante sea capaz de:

\begin{enumerate}
\item Construir atlas de variedades suaves y verificar la suavidad de transformaciones entre variedades.
\item Calcular espacios tangentes, diferenciales, fibrados tangente y cotangente, campos vectoriales y corchetes de Lie en ejemplos básicos, y describir los fibrados tangente y cotangente mediante sus trivializaciones locales.
\item Usar el lenguaje de fibrados, tensores y formas diferenciales para formular y resolver problemas de manera independiente de coordenadas.
\item Calcular pullbacks y derivadas exteriores, distinguir formas cerradas de formas exactas y explicar la idea básica de la cohomología de de~Rham.
\item Definir la integral de formas sobre variedades orientadas y aplicar el teorema de Stokes, reconociendo sus versiones clásicas.
\item Explicar las nociones de conexión afín, métrica riemanniana, conexión de Levi-Civita, geodésica y curvatura, efectuar cálculos en ejemplos básicos y relacionar la curvatura seccional con la curvatura gaussiana de superficies.
\item Consultar literatura especializada y comunicar argumentos de geometría diferencial con rigor, tanto de forma oral como escrita.
\end{enumerate}

\subsubsection*{Contenidos}

La distribución por semanas es orientativa (ciclo de aproximadamente 16~semanas). La persona docente puede reordenar subtemas y ajustar su profundidad, garantizando los contenidos básicos de cada bloque según los alcances indicados en la metodología.

\begin{enumerate}
\item \textbf{Variedades diferenciales y topología diferencial (5 semanas):}
\begin{enumerate}
    \item Variedades topológicas y suaves; cartas, atlas y funciones de transición.Condiciones de Hausdorff y segunda numerabilidad.
    \item Transformaciones suaves y difeomorfismos. Vectores tangentes y diferencial.
    \item Inmersiones, submersiones, embeddings y subvariedades. Valores regulares, teorema del valor regular e intersección transversal.
    \item Fibrados tangente y cotangente; trivializaciones locales.
    \item Nociones de fibrados vectoriales y secciones.
    \item Campos vectoriales, curvas integrales y flujos locales; corchete de Lie.
    \item Nociones elementales y ejemplos básicos de grupos de Lie y mapa exponencial.
\end{enumerate}

\item \textbf{Formas diferenciales (3 semanas):}
\begin{enumerate}
    \item Nociones de álgebra tensorial y exterior; tensores y formas diferenciales.
    \item Producto exterior y pullback.
    \item Derivada exterior, producto interior y derivada de Lie; fórmula de Cartan.
    \item Formas cerradas y exactas; lema de Poincaré; idea básica de la cohomología de de Rham y ejemplos.
\end{enumerate}

\item \textbf{Integración en variedades (3 semanas):}
\begin{enumerate}
    \item Orientación y variedades con frontera; orientación inducida en la frontera.
    \item Formas de soporte compacto y particiones de la unidad.
    \item Integración de formas de grado máximo y cambio de variables.
    \item Teorema de Stokes y recuperación de los teoremas fundamentales del cálculo vectorial.
\end{enumerate}

\item \textbf{Geometría riemanniana (4 semanas):}
\begin{enumerate}
    \item Conexiones afines, derivada covariante y transporte paralelo.
    \item Métricas riemannianas y longitud de curvas. Conexión de Levi-Civita: compatibilidad métrica y torsión nula. Geodésicas.
    \item Tensor de curvatura y curvatura seccional: definiciones y ejemplos básicos. Relación con la curvatura gaussiana de superficies e interpretación del teorema egregio.
\end{enumerate}

\end{enumerate}

\subsubsection*{Metodología}

\noindent \textbf{Lineamientos metodológicos:} La persona docente a cargo de este curso debe respetar las siguientes pautas:
\begin{enumerate}
    \item La presentación debe priorizar el rigor: definiciones precisas, enunciados completos y demostraciones de los resultados centrales. En particular, se demostrarán el teorema del valor regular, una versión local del lema de Poincaré y el teorema de Stokes. Los resultados auxiliares cuya prueba exceda el alcance del curso podrán utilizarse con enunciados e hipótesis precisos.
    \item Al estudiar campos vectoriales, curvas integrales y flujos, debe explicitarse el vínculo con la existencia y unicidad de soluciones de ecuaciones diferenciales ordinarias vistas en MA-0556. También debe enunciarse y utilizarse la dependencia suave respecto de los datos iniciales cuando sea necesaria.
    \item Cuando la naturaleza del contenido lo permita, se deben utilizar representaciones visuales que apoyen la comprensión de los temas.
    \item Se fomenta la comunicación clara y rigurosa de argumentos, tanto oral como escrita. Se incluirá al menos una actividad breve que permita practicar ambas formas de comunicación.
  \item La dedicación del estudiantado debe ser compatible con el creditaje. Los fibrados y el álgebra tensorial se desarrollarán en la medida necesaria para estudiar campos, formas y métricas; los grupos de Lie se introducirán mediante ejemplos y la noción de exponencial. Se garantizarán las definiciones y ejemplos básicos del bloque riemanniano, incluida la curvatura; las demostraciones y aplicaciones adicionales se ajustarán al tiempo disponible.
\end{enumerate}

\textbf{Sugerencias metodológicas:} Adicionalmente, se tienen las siguientes recomendaciones para la persona docente a cargo de este curso:
\begin{enumerate}
    \item Utilizar el recurso tecnológico para reforzar distintos conceptos estudiados en el curso por medio de la visualización, cálculo y experimentación.
    \item Aprovechar momentos adecuados del curso para motivar el estudio por medio de reseñas acerca del desarrollo histórico de la geometría diferencial y su importancia en aplicaciones dentro y fuera de la matemática.
    \item En la evaluación, se sugiere utilizar estrategias que promuevan el trabajo constante de parte de las personas estudiantes a lo largo del ciclo lectivo. Por ejemplo, esto se puede hacer por medio de tareas, quizzes o exposiciones breves sobre un tema o aplicación.
    \item Para fomentar el desarrollo de las habilidades investigativas en el estudiantado, se sugiere la implementación de pequeños proyectos de investigación. Por ejemplo, pueden ser proyectos de investigación bibliográfica o también proyectos cortos donde se desarrolle un tema por medio de ejercicios dirigidos.
    \item Dado que hoy en día la mayoría del trabajo en matemática, tanto a nivel académico como profesional, es de naturaleza colaborativa, se sugiere a la persona docente implementar estrategias que promuevan el trabajo en equipo.
    \item Explicitar al inicio del ciclo qué partes del bloque riemanniano se considerarán prioritarias, de modo que el estudiantado pueda planificar el trabajo independiente.
\end{enumerate}

\subsubsection*{Evaluación}

El curso se evaluará por medio de las siguientes actividades:

\begin{enumerate}
\item Exámenes.
\item Quices.
\item Tareas.
\item Análisis de artículos relacionados con el material visto en clase.
\item Proyectos.
\item Exposiciones.
\end{enumerate}

\nocite{Tu11}
\nocite{Lee13}
\nocite{Tu17}
\nocite{DoC92}
\nocite{GP74}
\nocite{DoC16}
\nocite{Kuh15}
\nocite{ONeil06}

\printcursosbibliography
\end{refsection}
